\begin{table}[htbp]
\centering
\caption{El sector energético en el artículo 1o., por vía de entrada (mdp)}
\label{tab:C3_1}
\small
\begin{tabular}{lSSSSS}
\toprule
{Concepto} & {2024} & {2025} & {Dif. nominal} & {Dif. real} & {Var. real \%} \\
\midrule
{Transferencias del Fondo Mexicano del Petroleo (ordinarias)} & \num{277774.3} & \num{279766.8} & \num{1992.5} & \num{-9951.8} & \num{-3.4} \\
{Ingreso propio de Petroleos Mexicanos} & \num{769805.6} & \num{860868.2} & \num{91062.6} & \num{57961.0} & \num{7.2} \\
{Ingreso propio de la Comision Federal de Electricidad} & \num{446951.3} & \num{539145.6} & \num{92194.3} & \num{72975.4} & \num{15.7} \\
{Impuesto por la actividad de exploracion y extraccion de hidrocarburos} & \num{7811.0} & \num{7140.1} & \num{-670.9} & \num{-1006.8} & \num{-12.4} \\
{SUMA de las cuatro vias} & \num{1502342.2} & \num{1686920.7} & \num{184578.5} & \num{119977.8} & \num{7.7} \\
\bottomrule
\end{tabular}
\par\vspace{2pt}\footnotesize Fuente: elaboración propia del ITED con la Ley de Ingresos aprobada de 2024 y 2025. Las diferencias en millones de pesos se reportan dos veces, nominal y real; la real está en pesos de 2025 con el deflactor de 4.3\% que declara el CGPE.
\end{table}
